\scanpage{161}
Si $2$ inv., la catégorie des modules à involution équivaut à celle des \deleted{\unsure{paire|paires}} couples de modules $(\underset{\substack{\shortparallel\\ E^0}}{E^-},E^+)$, et l'opération $*$ est l'opération de produit tensoriel séparément sur les deux facteurs \deleted{\unsure{linéaires|linéaire}} -- aucun mystère ! L'opération est délicate et utile, par contre, si $2$ non inversible. (p. ex. en car. $2$)

\begin{center}\rule{1.8em}{0.4pt}\end{center}

Revenons au contexte des espaces affines. Si on a un espace affine $A$ sous un module de translations \uline{inversible} \inserted{$L=A^0$}, alors la donnée d'un diviseur quadratique dans $\check{\mathbb V}(A)$ (le schéma \deleted{\unsure{espace}} \inserted{torseur sous $\check{\mathbb V}(A^0)$} défini par $A$) revient à la donnée d'une forme quadratique sur $E$ (enveloppe affine de $A$) à valeurs dans $L^{\otimes2}$
\[
q_E:E\longrightarrow L^{\otimes2}
\]
(i.e. d'une ``fonction \deleted{\unsure{pol.|polyn.}} \inserted{$q_A$} de degré $\leq2$ sur $A$''), telle que $q_L=q\mid L$ (la ``partie hom de $q$'') soit \deleted{\unsure{l'élem|l'élévation}} \deleted{forme} quadr. can
\[
L\longrightarrow L^{\otimes2}\qquad u\longmapsto u^{\otimes2}
\]
% archive sheet number: 160

